\section{Register overview}\label{sec:register-overview}

In this part, we implement the quantum registers.
Our goal is force Alice and Bob to share a state of the following form:
\newcommand{\attemptedwidth}{1cm}
\begin{center}
\begin{tabular}{|K{\attemptedwidth}|K{\attemptedwidth}|K{\attemptedwidth}|K{\attemptedwidth}|K{\attemptedwidth}|K{.5cm}|K{\attemptedwidth}|}
\cline{1-2}\cline{4-5}\cline{7-7}
$r_1$ & $r_2$ & $\cdots$ & $r_{k-1}$ & $r_k$ & $\otimes$ & $\mathrm{aux}$\\
\cline{1-2}\cline{4-5}\cline{7-7}
\end{tabular}~,
\end{center}
in which each register $r_i$ contains an EPR state,
and $\mathrm{aux}$ is a symmetric auxiliary state.
In addition, we want the verifier to be able to (i) force
the provers to perform Pauli basis queries on some of these registers and report back the outcomes
and (ii) ``hide" the remaining registers from the provers so that they do not measure them at all.

\subsection{Definitions}\label{sec:protocol-defs}

In this section,
we will begin by defining quantum registers
for \emph{nonuniform} games.
Defining registers for uniform games~$\game$ is a little more
complicated because we allow the number and size of registers for $\game(\mathsf{input})$
to depend on $\mathsf{input}$.
We detail this below in \Cref{sec:uniform-registers}.

\begin{definition}\label{def:register}
Let $k \geq 0$ be an integer, and let $n = (n_1, \ldots, n_k)$ and $q = (q_1, \ldots, q_k)$ be $k$-tuples of integers.
A \emph{$(k, n, q)$-register game} $\game$ is defined as follows.
\begin{itemize}
\itemsep -.5pt
\item[$\circ$] Questions~$x$ are formatted into two blocks $x = (x_1, x_2)$. The first block contains a list of~$k$ Pauli basis queries
			$x_1 = (W_1, \ldots, W_k)$, where each $W_i \in \{X, Z, \hideq, \noop\}$.
\item[$\circ$] Answers~$a$ are formatted into two blocks $a = (a_1, a_2)$. The first block contains a list of answers to the Pauli basis queries $a_1 = (u_1, \ldots, u_k)$.
			Here each $u_i  \in \F_{q_i}^{n_i} \cup \{\varnothing\}$.
\end{itemize}
An \emph{$(k, n, q)$-register strategy} $\calS$ is defined as follows.
\begin{itemize}
\itemsep -.5pt
\item[$\circ$] Alice and Bob share a state
\begin{equation*}
\ket{\psi} = \ket{r_1} \otimes \cdots \otimes \ket{r_k} \otimes \ket{\mathrm{aux}}.
\end{equation*}
Here, $\ket{r_i} = \ket{\mathrm{EPR}_{q_i}^{n_i}}$ for each~$i$, and $\ket{\mathrm{aux}}$ is an arbitrary symmetric shared state.
\item[$\circ$] Given a question $x = (x_1, x_2)$ with first block $x_1 = (W_1, \ldots, W_k)$, Alice and Bob act as follows.
	Let $i \in [k]$.
	\begin{itemize}
	\item If $W_i \in \{X, Z\}$, they measure $\tau^{W}$ on the $i$-th EPR register and set $u_i$ to be the outcome.
	\item If $W_i \in \{\hideq, \noop\}$, they set $u_i = \varnothing$.
	\end{itemize}
	Introduce the notation $\tau_{\varnothing}^W = I$ for $W \in \{\hideq, \noop\}$.
	We can write their measurement as
	\begin{equation}\label{eq:in-math}
		M^{x_1, x_2}_{a_1} = \tau^{W_1}_{u_1} \otimes \cdots \otimes \tau^{W_k}_{u_k} \otimes I_{\reg{aux}}.
	\end{equation}
	To produce the second part of their answer~$a_2$, Alice and Bob can measure any part of their state except the EPR registers which have been ``hidden".	
	This entails the following: let $S = \{i \mid W_i = \hideq\}$.  Then for any answer~$a$, the corresponding POVM acts as follows:
	\begin{equation}\label{eq:hide-coords-in-S}
	M_a^x = M_{\overline{S}} \otimes I_{S}.
	\end{equation}
	Here, $I_{S}$ is the identity matrix on the EPR registers in~$S$,
	whereas $M_{\overline{S}}$ is a POVM acting on the EPR registers in~$\overline{S}$ as well as the state~$\ket{\mathrm{aux}}$.
\end{itemize}
We define $\valreg{k,n,q}{\game}$ to be the maximum over $\valstrat{\game}{\calS}$, where $\calS$ is any $(k, n, q)$-register strategy.
\end{definition}

The~$X$ and~$Z$ questions specify the corresponding Pauli basis measurement,
and the~$\hideq$ question specifies that the register is to be hidden.
The~$\noop$ question is a ``no-op" and does not restrict Alice and Bob at all, other than making them respond with the ``no-op" answer $\varnothing$.
Thus, unlike with the data hiding question, they are allowed to measure the register as they see fit.
This will be useful later when we want Alice and Bob to measure both~$X$ \emph{and}~$Z$ observables on the same register.

In designing our compiler,
it will be convenient to define a set of strategies called ``semiregister strategies".
These will be strategies which are intermediate between $(k-1)$-register strategies and $k$-register strategies
in the sense that they have Pauli basis queries implemented on the final ($k$-th) register but not data hiding queries.
These are defined as follows.

\begin{definition}
A \emph{$(k, n, q)$-semiregister strategy} is defined just as a $(k, n, q)$-register strategy, with the following modification:
the set~$S$ used in \Cref{eq:hide-coords-in-S} is changed to be $S = \{i \neq k \mid W_i = \hideq\}$.
We define $\valsemi{k,n,q}{\game}$ to be the maximum over $\valstrat{\game}{\calS}$, where $\calS$ is any $(k, n, q)$-semiregister strategy.
\end{definition}
\noindent
Thus, querying the $k$-th register of a semiregister strategy with
a~$\hideq$ is the same as querying it with a~$\noop$.

The following lemma shows that we can restrict to \emph{projective}
register strategies without loss of generality.
\begin{lemma}\label{lem:proj-suffices}
  Let $\calS$ be a $(k, n,q)$-register strategy. Then there exists a
  $(k,n,q)$-register strategy $\calS'$ in which all measurements are
  projective, and which produces the same bipartite correlation as $\calS$.
\end{lemma}
\begin{proof}
  Start with the strategy $\calS$, and let the measurements be denoted
  $M^{x_1, x_2}_{a_1, a_2}$. From the definition of register
  strategies, we know that for every set of questions $x_1, x_2$, the
  corresponding measurement can be written as a product
  \[ M^{x_1, x_2}_{a_1,a_2} = (\tau^{W_{x_1}}_{a_1})_{S} \ot (A^{x_1, x_2, a_1}_{a_2})_{\ol{S}}, \]
  where $S$ is the set of registers which receive a Pauli basis query
  in the set $X, Z, H$, $\ol{S}$ is its complement, and the operators
  $\{A^{x_1, x_2, a_1}_{a_2}\}$  form valid POVMs with outcomes $a_2$
  for every choice of $x_1, x_2, a_1$.
  We will apply Naimark's theorem~\Cref{thm:naimark} using the
  universal auxiliary state $\ket{\rmaux}$ to the $A$ operator to produce projectors
  $A'^{x_1, x_2, a_1}_{a_2}$. Using these, we define a projective measurement
  \[ M'^{x_1, x_2}_{a_1, a_2} = (\tau^{W_{x_1}}_{a_1})_{S} \ot
    (A'^{x_1, x_2, a_1}_{a_2})_{\ol{S}}. \]
  It is not hard to see that $M'$ and $\ket{\rmaux}$ form a valid
  Naimark dilation of $M$.
  Let $\calS'$ be the strategy $\calS$ with the shared state
  $\ket{\psi}$ replaced by $\ket{\psi} \ot \ket{\rmaux_A} \ot
  \ket{\rmaux_B}$ and the measurements $M$ replaced by $M'$. By
  construction, $\calS'$ is a projective strategy. Further,
  from~\Cref{cor:bipartite-naimark}, it follows that the bipartite
  correlations produced by the strategies $\calS'$ and $\calS$ are the same.
\end{proof}
\subsection{Results}\label{sec:register-results}

The key elements of our compiler are two new nonlocal games
called the Pauli basis test and the data hiding game.
The Pauli basis test ensures that the provers share an EPR state and honestly answer Pauli basis queries to this state.
The data hiding game allows us to ``hide" this state from the provers, ensuring that they do not use this register unless we ask them to.

Our compiler operates a register at a time and involves two subroutines, $\calC_{k \rightarrow \mathrm{semi}}$ and $\calC_{\mathrm{semi}\rightarrow k-1}$.
Given a $k$-register game, $\calC_{k\rightarrow \mathrm{semi}}$ produces a $k$-semiregister game.
To do so, it removes the guarantee that the provers data hide the $k$-th register and replaces it by playing the data hiding game on this register.
Thus, although the provers are no longer forced to hide the $k$-th register, they will have to do so anyway if they want to pass the data hiding game.
Similarly, given a $k$-semiregister game, $\calC_{\mathrm{semi}\rightarrow k-1}$ produces a $(k-1)$-register game.
To do so, it removes the guarantee that the provers have a $k$-th EPR register and replaces it by playing the Pauli basis test.
Thus, by alternating these two subroutines, we can compile a $k$-register game into a $0$-register game, i.e.\ a general game.


Before giving the properties of the Pauli basis compiler, we will need two definitions.

\begin{definition}
Given a string $x = (x_1, \ldots, x_k)$ and an integer $0 \leq \ell \leq k$, write $\shorten{x}{\ell} := (x_1, \ldots, x_\ell)$.
We extend this to register parameters $\tau = (k, n, q)$ by setting 
$\shorten{\tau}{\ell} := (\ell, \shorten{n}{\ell}, \shorten{q}{\ell})$.
Thus, $\shorten{\tau}{\ell}$ is the register parameters for the first~$\ell$ registers of~$\tau$.
\end{definition}

\begin{definition}
Let~$n$ and~$q$ be integers and~$\eta$ be a real number. We say they \emph{satisfy the Pauli basis condition} if
\begin{equation*}
q = 2^t,\qquad
\frac{1}{\mathrm{\poly}(n)} \leq \eta \leq \frac{1}{2},
\qquad \frac{64 \log(n)^2}{\eta^2} \leq q \leq \poly(n).
\end{equation*}
\end{definition}

The following theorem describes the Pauli basis compiler.

\begin{theorem}\label{theorem:semi-to-k-1-layer}
Let $\register = (k, n, q)$, and let $n_k$, $q_k$, and $\eta$ satisfy the Pauli basis condition.
Suppose $\game_{\mathrm{semi}}$ is a $\register$-semiregister game,
and consider the $\shorten{\register}{k-1}$-register game $\game_{k-1} = \calC_{\mathrm{semi} \rightarrow (k-1)}(\game_{\mathrm{semi}})$.
\begin{itemize}
\item[$\circ$] \textbf{Completeness:} Suppose there is a value-$1$ $\register$-semiregister strategy for $\game_{\mathrm{semi}}$ which is also a real commuting EPR strategy.
	Then there is a value-$1$ $\shorten{\register}{k-1}$-register strategy for $\game_{k-1}$ which is also a real commuting EPR strategy.
\item[$\circ$] \textbf{Soundness:} If $\valreg{\shorten{\register}{k-1}}{\game_{k-1}}\geq 1-\eps$ then $\valsemi{\register}{\game_{\mathrm{semi}}} \geq 1 -\delta(\eps)$, where $\delta(\eps) = \mathrm{poly}(\eps, \eta)$.
\end{itemize}
Furthermore,
\begin{align*}
\qtime{\game_{k-1}} &= \qtime{\game_{\mathrm{semi}}} + O(\log(n_k)), \\
\qlength{\game_{k-1}} &= \qlength{\game_{\mathrm{semi}}} + O(\log(n_k)),  \\
\atime{\game_{k-1}} &= \atime{\game_{\mathrm{semi}}} + \poly(n_k),\\
\alength{\game_{k-1}} &= \alength{\game_{\mathrm{semi}}} + O(n_k \cdot \log \log(n_k)).
\end{align*}
\end{theorem}

The following theorem describes the data hiding compiler.

\begin{theorem}\label{theorem:k-to-semi-layer}
Suppose $\game_k$ is a $(k, n, q)$-register game,
and consider the $(k, n, q)$-semiregister game $\game_{\mathrm{semi}} = \calC_{k \rightarrow \mathrm{semi}}(\game_k)$.
\begin{itemize}
\item[$\circ$] \textbf{Completeness:} 
	Suppose there is a value-$1$ $(k,n,q)$-register strategy for $\game_{k}$ which is also a real commuting EPR strategy.
	Then there is a value-$1$ $(k,n,q)$-semiregister strategy for $\game_{\mathrm{semi}}$ which is also a real commuting EPR strategy.
\item[$\circ$] \textbf{Soundness:} If $\valsemi{k,n,q}{\game_{\mathrm{semi}}} \geq 1-\eps$ then $\valreg{k,n,q}{\game_{k}} \geq 1 -\delta(\eps)$, where $\delta(\eps) = \mathrm{poly}(\eps)$.
\end{itemize}
Furthermore,
\begin{equation*}
\qtime{\game_{\mathrm{semi}}} = O(\qtime{\game_k}), \quad
\atime{\game_{\mathrm{semi}}} = O(\atime{\game_k}),
\end{equation*}
\begin{equation*}
\qlength{\game_{\mathrm{semi}}} = O(\qlength{\game_k}), \quad
\alength{\game_{\mathrm{semi}}} = O(\alength{\game_k}).
\end{equation*}
\end{theorem}

Combining \Cref{theorem:semi-to-k-1-layer,theorem:k-to-semi-layer} gives us the main result of \Cref{part:stack},
a compiler~$\calC$ which compiles $k$-register games into general games.

\begin{theorem}\label{thm:registers}
Let $\game_k$ be a $(k, n, q)$-register game. 
Let $\eta = (\eta_1, \ldots, \eta_k)$, and suppose $n_i$, $q_i$, and $\eta_i$ pass the Pauli basis condition for all $i \in [k]$.
Write
\begin{equation*}
\game = \calC(\game_k) :=
\calC_{\mathrm{semi}\rightarrow 0}(\calC_{\mathrm{1} \rightarrow \mathrm{semi}}(~\cdots~
	\calC_{\mathrm{semi}\rightarrow k-1}(\calC_{\mathrm{k \rightarrow \mathrm{semi}}}(\game_k)))).
\end{equation*}
\begin{itemize}
\item[$\circ$] \textbf{Completeness:} Suppose there is a value-$1$ $(k,n,q)$-register strategy for $\game_{k}$ which is also a real commuting EPR strategy.
	Then there is a real commuting EPR strategy for~$\game$ with value~$1$.
\item[$\circ$] \textbf{Soundness:} If $\val{\game}\geq 1-\eps$ then $\valreg{(k,n,q)}{\game_{k}} \geq 1 -\delta(\eps)$, where $\delta(\eps) = \mathrm{poly}(\eps, \eta_1, \ldots, \eta_k)$.
\end{itemize}
Furthermore,
\begin{align*}
\qtime{\game} &= \qtime{\game_k} + O(\log(n_1))+ \cdots + O(\log(n_k)), \\
\qlength{\game} &= \qlength{\game_k} + O(\log(n_1)) + \cdots + O(\log(n_k)),  \\
\atime{\game} &= \atime{\game_k} + \poly(n_1) + \cdots + \poly(n_k),\\
\alength{\game} &= \alength{\game_k} + O(n_1 \cdot \log \log(n_1)) + \cdots O(n_k \cdot \log \log(n_k)).
\end{align*}
\end{theorem}

\subsection{Registers for uniform games}\label{sec:uniform-registers}

In this section, we generalize the notion of registers to the case of uniform games,
in which a different set of register parameters might be used for each input.
To compile these games, we will need for the register parameters themselves to be uniformly generated.

\begin{definition}\label{def:reg-params-generator}
Let $M_{\mathrm{Params}}$ be a Turing machine which,
given an input $\mathsf{input}$, outputs $\lambda = (k, n, q)$.
Let $\game$ be a (nonuniform) game.
Then we say $M_{\mathrm{Params}}$ \emph{outputs the register parameters of $\game$}
if for every $\mathsf{input}$, $\game(\mathsf{input})$ is a $M_{\mathrm{Params}}(\mathsf{input})$-register game.
\end{definition}

Given this, our compiler for uniform games is given as follows.

\begin{corollary}\label{thm:uniform-registers}
Let $\game(\cdot)$ be a (uniform) game,
and let $M_{\mathrm{Params}}$ be a Turing machine which outputs its register parameters.
Then there exists a (uniform) game $\game_{\mathrm{Compile}}(\cdot)$ with the following properties.
Given an input $\mathsf{input}$, write $\game:=\game(\mathsf{input})$,
$\game_{\mathrm{Compile}} := \game_{\mathrm{Compile}}(\mathsf{input})$,
and $\lambda = (k, n, q) := M_{\mathrm{Params}}(\mathsf{input})$.
\begin{itemize}
\item[$\circ$] \textbf{Completeness:} Suppose there is a value-$1$ $(k,n,q)$-register strategy for $\game$ which is also a real commuting EPR strategy.
	Then there is a real commuting EPR strategy for~$\game_{\mathrm{Compile}}$ with value~$1$.
\item[$\circ$] \textbf{Soundness:}
Let $\eta = (\eta_1, \ldots, \eta_k)$, and suppose $n_i$, $q_i$, and $\eta_i$ pass the Pauli basis condition for all $i \in [k]$.
If $\val{\game_{\mathrm{Compile}}}\geq 1-\eps$ then $\valreg{\lambda}{\game} \geq 1 -\delta(\eps)$, where $\delta(\eps) = \mathrm{poly}(\eps, \eta_1, \ldots, \eta_k)$.
\end{itemize}
Furthermore,
\begin{align*}
\qtime{\game} &= \qtime{\game_k} + O(\log(n_1))+ \cdots + O(\log(n_k)) + \mathsf{time}(M_{\mathrm{Params}}(\mathsf{input})), \\
\qlength{\game} &= \qlength{\game_k} + O(\log(n_1)) + \cdots + O(\log(n_k)),  \\
\atime{\game} &= \atime{\game_k} + \poly(n_1) + \cdots + \poly(n_k)+ \mathsf{time}(M_{\mathrm{Params}}(\mathsf{input})),\\
\alength{\game} &= \alength{\game_k} + O(n_1 \cdot \log \log(n_1)) + \cdots O(n_k \cdot \log \log(n_k)).
\end{align*}
\end{corollary}

\begin{proof}
We first compute $\lambda = M_{\mathrm{Params}}(\mathsf{input})$
in time $\mathsf{time}(M_{\mathrm{Params}}(\mathsf{input}))$.
Then it can be checked that the compiled game $\calC(\game(\mathsf{input}))$ from \Cref{thm:registers}
can be efficiently simulated given the register parameters~$\lambda$.
\end{proof}

\subsection{Organization}

The remainder of \Cref{part:stack} is organized as follows.
\begin{itemize}
\item In \Cref{sec:self-test-pauli}, we introduce the Pauli basis self-test and prove its correctness.
\item \Cref{sec:pauli-basis-compiler} implements the Pauli basis compiler.
\item In \Cref{sec:data-hiding-layer}, we introduce the data hiding game.
\item \Cref{sec:compile-sec} implements the data hiding compiler.
\item \Cref{sec:rotated-data-hiding} contains a generalization of the data hiding game which allows us to hide more general sets of Pauli observables.
		This is not needed to implement the quantum registers, but it will be needed in \Cref{part:neexp} when designing the $\neexp$ protocol.
\end{itemize}

%%% Local Variables:
%%% mode: latex
%%% TeX-master: "main"
%%% End:
%%%---------- close: register-overview

%%%%%%%%%%% jump to pauli-basis
%%%---------- open: pauli-basis
%!TEX root = main.tex

